\documentclass[handout]{beamer}
\mode<presentation>
\usepackage{xcolor}
\usepackage[toc]{multitoc}
\usepackage{color}
\usepackage{graphicx}
\graphicspath{{../buku_ajar/buku_ajar_logika_matematika/figures/}}
\usepackage{tikz}
\usetikzlibrary{shapes,arrows,positioning,calc}
\usepackage{listings}
\usepackage{lmodern}
\usepackage{amsmath}
\usetheme[secheader]{Boadilla}
\usecolortheme{whale}
\setbeamertemplate{caption}[numbered]
\setbeamertemplate{bibliography item}[text]
\setbeamertemplate{navigation symbols}{}

\newcommand{\logic}[1]{\textbf{\textit{#1}}}

\theoremstyle{definition}
\newtheorem{definisi}[theorem]{Definisi}
\theoremstyle{example}
\newtheorem{contoh}[theorem]{Contoh}

\renewcommand*{\multicolumntoc}{2}

\setbeamertemplate{footline}
{
  \leavevmode%
  \hbox{%
  \begin{beamercolorbox}[wd=.333333\paperwidth,ht=2.25ex,dp=1ex,center]{author in head/foot}%
    \usebeamerfont{author in head/foot}\insertshortauthor%
  \end{beamercolorbox}%
  \begin{beamercolorbox}[wd=.433333\paperwidth,ht=2.25ex,dp=1ex,center]{title in head/foot}%
    \usebeamerfont{title in head/foot}\insertshorttitle
  \end{beamercolorbox}%
  \begin{beamercolorbox}[wd=.233333\paperwidth,ht=2.25ex,dp=1ex,right]{date in head/foot}%
    \usebeamerfont{date in head/foot}\insertshortdate{}\hspace*{2em} \insertframenumber{} / \inserttotalframenumber\hspace*{2ex} 
  \end{beamercolorbox}}%
  \vskip0pt%
}

\subtitle[Logika Matematika]{Logika Matematika}
\title[Kuantor]{\texorpdfstring{Kuantor Universal ($\forall$) dan Eksistensial ($\exists$)}{Kuantor Universal dan Eksistensial}}
\author{Cecep Suwanda, S.Si., M.Kom.}
\date{Maret 2025}
\institute{Universitas Bale Bandung}

\begin{document}

\maketitle

\section*{Daftar Isi}
\begin{frame}
\frametitle{Daftar Isi}
\tableofcontents
\end{frame}

% ============================================================
% SECTION 7-1: Pengantar Kuantifikasi
% ============================================================
\section{Pengantar Kuantifikasi}

\begin{frame}
\frametitle{Kuantifikasi}
\begin{definisi}[Kuantifikasi]
\logic{Kuantifikasi} adalah proses pemberian kuantor pada fungsi proposisional untuk menyatakan jangkauan objek. Operator disebut \logic{kuantor}.
\end{definisi}
Melalui kuantifikasi, pernyataan dapat menyatakan sifat untuk \textbf{seluruh} atau \textbf{sebagian} elemen domain.
\end{frame}

\begin{frame}
\frametitle{Semesta Pembicaraan ($\mathcal{U}$)}
\begin{definisi}[Semesta Pembicaraan]
Himpunan semua nilai yang diperbolehkan untuk menggantikan variabel pada fungsi proposisional.
\end{definisi}
\begin{contoh}
``Semua elemen $x$ memenuhi $x^2 > 0$''
\begin{itemize}
  \item $\mathcal{U} = \mathbb{R}$: \textbf{Salah} (counterexample: $x=0$)
  \item $\mathcal{U} = \mathbb{Z}^+$: \textbf{Benar}
\end{itemize}
\end{contoh}
\end{frame}

\begin{frame}
\frametitle{Dua Kuantor Utama}
\begin{enumerate}
  \item \textbf{Kuantor universal} ($\forall$): properti berlaku bagi \textbf{seluruh} elemen domain.
  \item \textbf{Kuantor eksistensial} ($\exists$): ada \textbf{sekurang-kurangnya satu} elemen yang memenuhi properti.
\end{enumerate}
\end{frame}

% ============================================================
% SECTION 7-2: Kuantor Universal
% ============================================================
\section{Kuantor Universal}

\begin{frame}
\frametitle{Definisi Kuantor Universal}
$\forall x \, P(x)$ = ``$P(x)$ benar untuk \textbf{semua} $x$ di $\mathcal{U}$''
\begin{itemize}
  \item Dibaca: ``Untuk setiap $x$, $P(x)$'' atau ``Untuk semua $x$, berlaku $P(x)$''
\end{itemize}
\begin{block}{Domain Terhingga}
$\mathcal{U} = \{a_1, \ldots, a_n\}$ $\Rightarrow$ $\forall x \, P(x) \equiv P(a_1) \land P(a_2) \land \ldots \land P(a_n)$
\end{block}
Satu \textit{counterexample} $P(c)$ = Salah $\Rightarrow$ klaim universal salah.
\end{frame}

\begin{frame}
\frametitle{Universal + Implikasi}
Pola ``Setiap A adalah B'' $\Rightarrow$ $\forall x \, (M(x) \rightarrow P(x))$
\begin{contoh}
``Setiap mahasiswa mematuhi peraturan''\\
$\mathcal{U}$ = mahasiswa: $\forall x \, P(x)$\\
$\mathcal{U}$ = makhluk hidup: $\forall x \, (M(x) \rightarrow P(x))$
\end{contoh}
\begin{alertblock}{Perhatian}
Jangan gunakan $\forall x \, (M(x) \land P(x))$ untuk ``setiap A adalah B''. Itu berarti seluruh semesta adalah A dan B.
\end{alertblock}
\end{frame}

\begin{frame}
\frametitle{Contoh Matematis Universal}
$\mathcal{U} = \mathbb{R}$:
\begin{itemize}
  \item $\forall x \, (x^2 \geq 0)$: \textbf{Benar}
  \item $\forall x \, (x + 1 > x)$: \textbf{Benar}
  \item $\forall x \, (x^2 > x)$: \textbf{Salah} (counterexample: $x = 0.5$)
\end{itemize}
\end{frame}

% ============================================================
% SECTION 7-3: Kuantor Eksistensial
% ============================================================
\section{Kuantor Eksistensial}

\begin{frame}
\frametitle{Definisi Kuantor Eksistensial}
$\exists x \, P(x)$ = ``Terdapat \textbf{setidaknya satu} $x$ di $\mathcal{U}$ sehingga $P(x)$ benar''
\begin{itemize}
  \item Dibaca: ``Ada $x$ dengan sifat $P(x)$'' atau ``Sekurang-kurangnya satu $x$ memenuhi $P(x)$''
\end{itemize}
\begin{block}{Domain Terhingga}
$\exists x \, P(x) \equiv P(a_1) \lor P(a_2) \lor \ldots \lor P(a_n)$
\end{block}
Satu \textit{witness} $P(c)$ = Benar $\Rightarrow$ klaim eksistensial benar.
\end{frame}

\begin{frame}
\frametitle{Eksistensial + Konjungsi}
Pola ``Sebagian A adalah B'' $\Rightarrow$ $\exists x \, (A(x) \land B(x))$
\begin{contoh}
``Sebagian buah terasa asam''\\
$\mathcal{U}$ = buah: $\exists x \, A(x)$\\
$\mathcal{U}$ = benda di dapur: $\exists x \, (B(x) \land A(x))$
\end{contoh}
\begin{alertblock}{Perhatian}
Jangan gunakan $\exists x \, (B(x) \rightarrow A(x))$---\textit{vacuous truth}: elemen bukan $B$ sudah cukup membuat implikasi benar.
\end{alertblock}
\end{frame}

\begin{frame}
\frametitle{Contoh Matematis Eksistensial}
$\mathcal{U} = \mathbb{R}$:
\begin{itemize}
  \item $\exists x \, (x^2 = 9)$: \textbf{Benar} (witness: $x=3$)
  \item $\exists x \, (x = x + 1)$: \textbf{Salah}
  \item $\exists x \, (x^2 < 0)$: \textbf{Salah} di $\mathbb{R}$; Benar di $\mathbb{C}$ ($x=i$)
\end{itemize}
\end{frame}

% ============================================================
% SECTION 7-4: Kuantor Bersarang
% ============================================================
\section{Kuantor Bersarang}

\begin{frame}
\frametitle{Definisi Kuantor Bersarang}
\begin{definisi}[Kuantor Bersarang]
Kuantor diletakkan di dalam \textit{scope} kuantor lain. Penulisan ditumpuk dari luar (kiri) ke dalam (kanan).
\end{definisi}
$\forall x \, \exists y \, P(x,y)$: untuk setiap $x$, terdapat $y$ (bergantung pada $x$) sehingga $P(x,y)$ benar.
\end{frame}

\begin{frame}
\frametitle{Empat Konfigurasi Bersarang}
\begin{center}
\small
\begin{tabular}{|l|l|}
\hline
\textbf{Ekspresi} & \textbf{Makna} \\ \hline
$\forall x \forall y \, P(x,y)$ & $P(x,y)$ benar untuk semua pasangan $(x,y)$ \\ \hline
$\forall x \exists y \, P(x,y)$ & Untuk setiap $x$, ada $y$ (bergantung $x$) \\ \hline
$\exists x \forall y \, P(x,y)$ & Ada satu $x$ sehingga $P(x,y)$ benar untuk semua $y$ \\ \hline
$\exists x \exists y \, P(x,y)$ & Ada minimal satu pasangan $(x,y)$ \\ \hline
\end{tabular}
\end{center}
\end{frame}

\begin{frame}
\frametitle{Contoh Matematis Bersarang}
$\mathcal{U} = \mathbb{R}$:
\begin{itemize}
  \item $\forall x \forall y \, (x+y = y+x)$: \textbf{Benar} (komutatif)
  \item $\forall x \exists y \, (x+y = 0)$: \textbf{Benar} (pilih $y = -x$)
\end{itemize}
\end{frame}

% ============================================================
% SECTION 7-5: Urutan Kuantor
% ============================================================
\section{Urutan Kuantor}

\begin{frame}
\frametitle{Asimetri Urutan}
$L(x,y)$ = ``$x$ mencintai $y$'', $\mathcal{U}$ = manusia
\begin{block}{$\forall x \, \exists y \, L(x,y)$}
Setiap orang punya minimal satu orang yang dicintai. Nilai $y$ dapat berbeda untuk setiap $x$.
\end{block}
\begin{block}{$\exists y \, \forall x \, L(x,y)$}
Ada satu orang yang dicintai oleh \textbf{semua} orang. Klaim lebih kuat!
\end{block}
\end{frame}

\begin{frame}
\frametitle{Kaidah Implikasi Urutan}
\[ \exists y \, \forall x \, P(x,y) \implies \forall x \, \exists y \, P(x,y) \]
Jika ada satu $y$ untuk semua $x$, maka setiap $x$ punya minimal satu $y$.

\begin{block}{Sebaliknya tidak berlaku}
``Setiap mahasiswa punya pacar'' $\not\Rightarrow$ ``Ada satu orang berpacaran dengan seluruh mahasiswa''
\end{block}
\end{frame}

% ============================================================
% SECTION 7-6: Negasi dan De Morgan
% ============================================================
\section{Negasi Berkuantor}

\begin{frame}
\frametitle{Hukum De Morgan untuk Kuantor}
\begin{block}{Negasi Universal}
$\neg \forall x \, P(x) \equiv \exists x \, \neg P(x)$\\
``Bukan semua pandai'' = ``Ada yang tidak pandai''
\end{block}
\begin{block}{Negasi Eksistensial}
$\neg \exists x \, P(x) \equiv \forall x \, \neg P(x)$\\
``Tidak ada yang tegas berlebihan'' = ``Semua tidak tegas berlebihan''
\end{block}
\end{frame}

\begin{frame}
\frametitle{Negasi Bersarang}
Setiap kuantor dilalui: $\forall \leftrightarrow \exists$
\begin{align*}
\neg \big[ \exists x \, \forall y \, \forall z \, P(x,y,z) \big] &\equiv \forall x \, \exists y \, \exists z \, \neg P(x,y,z)
\end{align*}
\end{frame}

\begin{frame}
\frametitle{Negasi dengan Implikasi}
``Semua mahasiswa baru mengikuti orientasi'': $\forall x \, (M(x) \rightarrow O(x))$
\begin{align*}
\neg\big(\forall x \, (M \rightarrow O)\big) &\equiv \exists x \, (M(x) \land \neg O(x))
\end{align*}
\textbf{Makna:} ``Terdapat mahasiswa baru yang tidak mengikuti orientasi.''
\end{frame}

% ============================================================
% SECTION 7-7: Panduan Translasi
% ============================================================
\section{Panduan Translasi}

\begin{frame}
\frametitle{Kamus Translasi (1/2)}
\begin{center}
\small
\begin{tabular}{|l|l|}
\hline
\textbf{Frasa} & \textbf{Formal} \\ \hline
Setiap $P$ itu $Q$ & $\forall x \, (P(x) \rightarrow Q(x))$ \\ \hline
Ada/Beberapa $P$ yang $Q$ & $\exists x \, (P(x) \land Q(x))$ \\ \hline
Tidak ada $P$ yang $Q$ & $\forall x \, (P(x) \rightarrow \neg Q(x))$ \\ \hline
Beberapa $P$ TIDAK $Q$ & $\exists x \, (P(x) \land \neg Q(x))$ \\ \hline
\end{tabular}
\end{center}
\end{frame}

\begin{frame}
\frametitle{Contoh Translasi}
$\mathcal{U}$ = makhluk. $A(x)$=atlet, $S(x)$=sehat, $K(x)$=perokok.
\begin{itemize}
  \item Setiap atlet pasti sehat: $\forall x \, (A(x) \rightarrow S(x))$
  \item Tidak ada perokok yang atlet: $\forall x \, (K(x) \rightarrow \neg A(x))$
  \item Beberapa atlet sehat juga perokok: $\exists x \, (A(x) \land S(x) \land K(x))$
  \item Sebagian atlet tidak sehat: $\exists x \, (A(x) \land \neg S(x))$
  \item Hanya yang tidak merokok yang sehat: $\forall x \, (S(x) \rightarrow \neg K(x))$
\end{itemize}
\end{frame}

% ============================================================
% RANGKUMAN
% ============================================================
\section{Rangkuman}

\begin{frame}
\frametitle{Rangkuman}
\begin{itemize}
  \item $\mathcal{U}$ menentukan domain; memengaruhi nilai kebenaran.
  \item $\forall$ + implikasi untuk ``setiap A adalah B''; $\exists$ + konjungsi untuk ``ada A yang B''.
  \item Urutan kuantor penting: $\forall x \exists y \neq \exists y \forall x$.
  \item De Morgan: $\neg\forall \equiv \exists\neg$; $\neg\exists \equiv \forall\neg$.
\end{itemize}
\end{frame}

\begin{frame}
\frametitle{Kata Kunci}
\begin{center}
semesta pembicaraan $\bullet$ kuantor universal $\bullet$ kuantor eksistensial\\
implikasi $\bullet$ konjungsi $\bullet$ kuantor bersarang $\bullet$ negasi kuantor\\
witness $\bullet$ counterexample
\end{center}
\end{frame}

\end{document}
